% Pursuing Stacks (cotes 134-2 to 134-8): archive cote 134-6, site page 55 (PDF page 56), typescript page 373.
% typescript. Draft compared: ps5.tex:4097-4159.
% AI-transcribed by Claude Opus 5.5 (claude-opus-5-5, adaptive thinking, effort high), Message Batches API, batch ps3 (msgbatch_01XdFmuVVBAcgc7EFxdtCqsj), 2026-11-12; not yet checked by eye.
\scanpage{134-6/55}
\mspage{373}

the writing of these notes. (Namely, first in connection with my reflec-%
tions on derivators\inserted{(cf. section 69)}, and now in connection with the reflections on abeli-%
anization.) In view of my reflections on derivators, I would like to
view the functor (18) as an ope\unsure{ra}tion of ``integration'', whereas the
traditional functor
\[ \text{(19)} \qquad Rf_{*} : D^{+}(B'^{\wedge}_{ab}) \hand{\longrightarrow} D^{+}(B^{\wedge}_{ab}) \]
is \deleted{a functor of}\inserted{viewed as} ``cointegration'' (which I prefer to my former way of cal-%
ling it an ``integration''). T\over{\unsure{h}}{h}e first should be viewed as expressing \uline{homo-}%
\uline{logy} properties of the map f (in (Cat)\unsure{\deleted{,}} \hand{(}or between the corresponding
topoi), just as the latter expresses \uline{cohomology} properties of f. \unsure{\deleted{(}}This
does check with the corresponding qualifications ``integration'' - ``cointe-%
gration'' - as well as with the intuition, when B is reduced to a point,
identifying the first to a kind of direct sum (= integration), whereas
the latter is viewed as a kind of direct product (= \uline{co}integration). The
idea behind the terminology \inserted{will} go\deleted{es} through maybe when looking at the parti-%
cular case when B' is \xout{axsxx} a sum of copies of B\hand{,} namely a product of B
by a disc\over{\unsure{c}}{r}ete category I , and \xout{f}
\[ f\colon B' = BxI \hand{\longrightarrow} B \]
the projection.

T\unsure{H}e point I want to make here, mainly to myself, is that in the
present context \uline{when} (18) \uline{namely integration exists}, \uline{this operation}\inserted{presumably} \uline{is}
\uline{by no means less meaningful and important than the familiar} $Rf_{*}$ or \uline{cointe-}%
\uline{gration} - or \uline{equivalently stated}, \uline{that the homology properties of f are}
\uline{just as meaningful and deserving close attention}, \deleted{than}\inserted{as} \uline{the cohomology}
\uline{properties}, which so far have been the only ones \unsure{\hand{I}} have been looking at\unsure{\hand{.}}
Presumably when following this recommandation, a few unexpected \xout{facts}
facts and relationships shou\over{\unsure{d}}{l}d come out, such as various ``duality'' rela-%
tionships between homology properties of f , and cohomology properties
of the map $f^{o}$ between the opposite categories. (This is suggested by some
of the scratchwork I made on derivators and cohomology properties of maps
in (Cat) .) The only trouble is that such change or broadening of empha-%
sis as I am now suggesting will require a certain amount of extra attenti-%
on, which I am not too sure to be willing to invest in the subject,
namely algebraic topology. Thus presumably my main emphasis will remain
\margin{\unsure{$\updownarrow$}} with cohomology, rather than homology. \inserted{I am no longer convinced though that this point of view is technically more adequate than the dual one.}

\hand{2)} All the re\unsure{fl}exions of yesterday's notes as well as today's can
be extended, when replacing throughout abelian \inserted{pre}\deleted{s}heaves by presheaves of

\margin{55}
